% Section: Methods — P8 commit 3.
% Numerical claims verified against:
%   - ODE_PDE_PRE_REG.md §§ 1-7
%   - PRE_REG_AMENDMENT.md A1--A11
%   - src/quantum_liquid_neuralode/evaluation/h1_verdict.py
%   - src/qlnn_/training/{pde_residual_loss, multi_state_solver,
%       classical_pinn_solver, p3_8_review_demo, p3_9_pde_matrix}.py
%   - Per-cell metrics.json (param counts)

\section{Methods}
\label{sec:methods}

\subsection{Pre-registration and the falsifiable hypothesis}
\label{sec:methods-prereg}

The complete pre-registration document (\texttt{ODE\_PDE\_PRE\_REG.md},
committed 2026-05-19 before any benchmark results were generated)
states the central hypothesis:

\begin{quote}
\textbf{H1.} Quantum PINNs achieve a regime-dependent advantage over
matched classical baselines: $\dsmooth - \dbroad > 0$ with
paired-bootstrap 95\% confidence interval excluding zero, where
$\Delta_\text{cell} \equiv \relltwo_\text{classical} -
\relltwo_\text{QLNN}$ and the mean is taken per regime.
\end{quote}

The theoretical motivation is the Schuld--Sweke--Meyer
result~\cite{schuld2021effect} that a data-reuploading parameterized
quantum circuit realizes a truncated Fourier series in the encoded
variable, with accessible $K_\text{max}$ proportional to reuploading
depth. A model whose hypothesis class is a low-order Fourier series
should exhibit an inductive bias toward smooth periodic dynamics and
a corresponding difficulty representing broadband or multiscale
structure. H1 operationalizes this prediction as a comparison
between two pre-registered regime bins.

The pre-registration locks three additional commitments before
any result is generated: (i)~mandatory classical baselines
(capacity-matched cPINN, non-liquid Neural-ODE, plain MLP
forecaster, and a known-structure skyline); (ii)~banned headline
metrics that admit persistence-trivial wins (notably one-step MAE
on autoregressive rollout); (iii)~explicit decision rules under
which H1 is declared CONFIRMED, FALSIFIED, or INCONCLUSIVE.

Nineteen amendments to the pre-registration are documented openly
in \texttt{PRE\_REG\_AMENDMENT.md} (A1--A19). They were necessary
because the original pre-registration left some parameter choices
implicit; each amendment locks a single decision and reports
whether the outcome was stable under the alternative choice. The
key amendments from the original analysis pass are: A1 (skyline
threshold = 0.5; sensitivity at 0.75 reported); A4 (MLP capacity
3.3$\times$ Neural-ODE, violating the original ``within factor of
2'' rule, disclosed and shown not to affect the verdict sign); A7
(strict-versus-raw underfit-guard reporting); A9 (symmetric
quantum HPO sensitivity); A10 (combined ODE+PDE verdict at
$n=18$); A11 (full-ladder expansion to $n=24$ + KdV/Kuramoto
deferral rationale). Amendments A15--A19 were filed 2026-05-28
during an internal audit that surfaced (a) cross-side
training-budget parity across all solver models, (b) a PennyLane
\texttt{StronglyEntanglingLayers} fallback that silently aliased
\texttt{data\_reuploading} at the project's qubit count, (c) a
\texttt{qcpinn} quantum-attribution sub-experiment with three
step-wise variants, (d) \texttt{brickwall}'s structural
connectivity deficit at small qubit count (removed from the
empirical sweep; T3 mechanism scalars retained as
untrained-circuit diagnostic data), and (e) cross-task
forecaster/solver training-step budget parity. The verdict
refresh under the A15--A19 configuration is deferred to a
supplementary compute campaign; the verdicts reported in this
manuscript reflect the configuration locked at the time of
analysis.

\subsection{Hardness ladder}
\label{sec:methods-ladder}

We test H1 on eight pre-registered dynamical systems organized
into a regime-tagged hardness ladder.

\textbf{Smooth/periodic bin} (six pre-reg cells, four taxa, two
mathematical structures):
\begin{itemize}
  \item \lv{} ($\alpha = 1.1, \beta = 0.4, \delta = 0.1,
        \gamma = 0.4$, $T = 5$) -- 2D predator--prey limit cycle.
  \item \vdp{} ($\mu = 5$, $T = 10$) -- 2D stiff but periodic
        relaxation oscillator; pre-registered as borderline
        smooth.
  \item Heat equation ($\nu = 0.1$, $T = 1$, periodic
        on $[0, 2\pi)$) -- analytic reference
        $u(t,x) = e^{-\nu t}\sin x$.
  \item Burgers smooth ($\nu = 0.12$, $T = 2$, periodic on
        $[0, 2\pi)$, IC $\sin x$) -- npz-reference field.
\end{itemize}

\textbf{Broadband/multiscale bin} (six pre-reg cells, four taxa,
two mathematical structures):
\begin{itemize}
  \item Lorenz ($\sigma = 10, \rho = 28, \beta = 8/3$, $T = 2$)
        -- 3D chaotic flow.
  \item \fhn{} ($I = 0.5, \varepsilon = 0.08, a = 0.7, b = 0.8$,
        $T = 20$) -- 2D fast--slow stiff relaxation.
  \item \ac{} ($\varepsilon = 0.06$, $T = 8$, periodic on
        $[0, 2\pi)$) -- sharp two-front equilibrium.
  \item Burgers shock ($\nu = 0.004$, $T = 2$, periodic on
        $[0, 2\pi)$, IC $\sin x$) -- near-shock at $t^* \approx 1$.
\end{itemize}

Two pre-registered systems are mechanism-proven-but-compute-deferred
(amendment A11): the 12-dimensional Kuramoto coupled-oscillator
system (per-component scalar circuits scale linearly in state
dimension, projecting $\approx 7$~hr per cell) and KdV
($u_t + 6 u u_x + u_{xxx} = 0$). The triple-nested reverse-mode
autodiff required for the KdV solver was gate-tested
(\texttt{scripts/run\_p7\_8\_kdv\_gate.py}, result:
PASS, $u_{xxx}$ values finite and nontrivial, JIT compile
4.55~s, per-point cost ratio 0.431$\times$ the second-order
baseline thanks to XLA fusion); the integrated training cost at
$32 \times 32$ collocation and 2400 steps remains prohibitive
($\approx 12$~s per gradient step, projecting $\approx 8$~hr
per seed across all four quantum families plus the classical
baseline). Both deferrals are queued for the follow-up paper.

\subsection{Quantum-PINN family roster}
\label{sec:methods-roster}

Four quantum-PINN ansatz families are exercised at all 24 cells.
Each family is implemented faithfully against its published
specification; the binding manifest is
\texttt{refs/CIRCUIT\_SPECS.md}. For solver-task ODEs, each $d$-component
state is fitted via $d$ independent scalar circuits with $d$
separate weight pytrees (no entanglement across components); for
PDEs, each family has a 2D port that encodes $(t, x)$ on split
qubit registers.

\begin{itemize}
  \item \textbf{Data-reuploading} (Pérez-Salinas et
        al.~\cite{perezsalinas2020dataru}): the universal
        classifier ansatz with interleaved angle-encoding and
        parameterized $R(\theta)$ blocks plus ring CNOT
        entanglers. 4 qubits, 5 layers, 60 PQC parameters per
        scalar circuit.
  \item \textbf{TE-QPINN/FNN} (Berger et
        al.~\cite{berger2025teqpinn}): trainable-embedding QPINN
        with classical FNN encoder feeding angle-encoded rotations.
        60 PQC parameters and 100 classical FNN parameters per
        scalar circuit.
  \item \textbf{TE-QPINN/QNN} (Berger et
        al.~\cite{berger2025teqpinn}): trainable-embedding QPINN
        with QNN encoder (zero classical parameters by
        construction). 84 PQC parameters per scalar circuit.
  \item \textbf{QCPINN} (Farea et
        al.~\cite{farea2025qcpinn}): hybrid quantum--classical
        PINN with classical pre- and post-net. 15 PQC parameters
        and 706 classical parameters per scalar circuit (the
        classical-parameter overhead is disclosed openly and
        flagged at every cell in the per-cell tables in
        Section~\ref{sec:solver}).
\end{itemize}

For all four families the 2D PDE port uses split qubit registers:
$n_t = n_x = 4$ qubits per coordinate axis, 5 hardware-efficient
layers, ring entanglement. The Chebyshev coordinate feature map
realizes a tower of Chebyshev polynomials in each axis variable;
the convention is \texttt{src/qlnn\_/training/pde\_residual\_loss.py}
\texttt{ChebyshevDQC2DConfig}.

\subsection{Matched classical baselines}
\label{sec:methods-baselines}

Per pre-registration §6 (mandatory), four classical baselines are
trained at every applicable cell at matched parameter count
within a factor of two (except A4-disclosed MLP):

\begin{itemize}
  \item \textbf{Classical PINN} (cPINN): a small MLP trained with
        the same Lagaris hard-IC trial solution
        $u(t, x) = u_0(x) + (t - t_0) \cdot \text{MLP}_\theta(t, x)$
        and the same physics residual as each quantum-PINN cell.
        For vector ODE: $u(t) = u_0 + (t - t_0) \cdot \text{MLP}(t)$
        with output dimension $d$. The classical solver-task control.
  \item \textbf{Plain non-liquid Neural-ODE}: an MLP cell
        $h' = W_{\text{out}}\tanh(W_h [h, x, t])$ with no learnable
        time constants, Diffrax integration, the same residual head
        as the quantum forecaster. The forecaster-task control.
  \item \textbf{Plain MLP forecaster}: $y = \text{MLP}(\text{history}, x)$
        as a capacity-matched non-recurrent forecaster control.
  \item \textbf{Known-structure skyline}: per pre-reg §6, fits only
        free coefficients (e.g.~$\nu$, $\rho$) of the ground-truth
        equation by least squares against the reference trajectory.
        Used as the underfit/out-of-reach gate.
\end{itemize}

\subsection{Training protocol}
\label{sec:methods-train}

All solver-task cells use the Lagaris hard-IC trial solution. The
ODE form is $u(t) = u_0 + (t - t_0) \cdot (s \cdot
\text{circuit}(t) + b)$ with $\{s, b\}$ scalar training parameters
per component; the PDE form is
$u(t, x) = u_0(x) + (t - t_0) \cdot
(s \cdot \text{circuit}(t, x) + b)$. The physics-residual loss is
the mean squared residual over interior collocation points;
all four quantum-PINN families and the classical PINN share the
same loss closure, so the comparison differs only in the function
class of the trial solution.

Per-axis collocation grids: 60 interior points for ODE solvers, and
either $24{\times}24$, $28{\times}28$, or $32{\times}32{-}64$
for PDE solvers (the AC and burgers\_shock grids are intentionally
finer to capture the sharper structures, per A11). Training step
budgets per system are documented in
\texttt{CORRECTED\_PDE\_CONFIGS}: 1200 for heat, 1500 for
burgers\_smooth, 1800 for AC, 2400 for burgers\_shock. ODE step
budgets follow \texttt{multi\_state\_solver.\_SCALAR\_FAMILIES}:
1200 (chebyshev\_dqc), 1500 (te\_qpinn\_fnn, qcpinn), 2000
(te\_qpinn\_qnn).

\subsection{Verdict machinery}
\label{sec:methods-verdict}

Per cell the trained model is evaluated on an interior 50$\times$50
grid (PDE) or a 1024-point time grid (ODE), and we record
relative-$L^2$ error against the reference field. We construct one
\texttt{CellRecord} per (system, seed) with the best quantum-PINN
relative-$L^2$ across the four families on the quantum side and
the matched classical-PINN result on the classical side. The H1
verdict module (\texttt{h1\_verdict.py} in the
\texttt{quantum\_liquid\_neuralode.evaluation} package) performs
the following:

\begin{enumerate}
  \item \emph{Underfit guard}: drops cells where the classical
        baseline fails to reach an adequacy threshold (its train-side
        $\relltwo > 0.5$ under A7's strict reading).
  \item \emph{Skyline guard}: drops cells where the known-structure
        skyline cannot itself reach the adequacy threshold.
        At threshold 0.5 (A1), all skyline cells pass.
  \item \emph{Paired-bootstrap}: resamples the kept cells with
        replacement, computing per-resample $\dsmooth$ and
        $\dbroad$, and reports the percentile 95\% CI of
        $\ddiff = \dsmooth - \dbroad$. Default $n_{\text{iter}} =
        10{,}000$, $\alpha = 0.05$.
  \item \emph{Decision rule}: CONFIRMED if the CI excludes zero
        positively; FALSIFIED if it includes zero or excludes
        zero negatively; INCONCLUSIVE if the guards drop too
        many cells to render the bootstrap underpowered.
\end{enumerate}

The pre-registration commits the framework to mechanical
application of these rules; no analyst chooses the verdict
post-hoc.

\subsection{Symmetric HPO sensitivity}
\label{sec:methods-hpo}

Amendment A9 adds a symmetric quantum-side hyperparameter sweep
that mirrors the classical-side sweep already in
\texttt{run\_p7\_5\_hpo\_sensitivity.py}. At three anchor cells
(\lv{} seed~2, \vdp{} seed~1, Lorenz seed~2), each
of the four quantum families is retrained at three learning
rates $\{10^{-3}, 5 \cdot 10^{-3}, 10^{-2}\}$ and two training
budgets $\{1500, 3000\}$ steps -- 72 retrains in total. The
``full HPO-best'' verdict reported in Section~\ref{sec:verdict}
takes the per-cell minimum-$\relltwo$ HPO point on both sides.
This is the Bowles--Schuld~\cite{bowles2024better} ``tune both
sides'' sensitivity test.

\subsection{Open-source reproducibility}
\label{sec:methods-repro}

The repository at
\url{https://github.com/shawngibford/qlnn} contains the complete
code, pre-registration, and committed result JSONs.
\texttt{scripts/reproduce\_paper.sh} regenerates every committed
result from scratch (estimated wall-clock $\approx$~8~hr on a
single MacBook Pro M1, parallelizable);
\texttt{scripts/verify\_paper\_integrity.py} mechanically checks
every cited number in this paper against its committed JSON file
and returns non-zero if any number drifts. Tagged GitHub releases
mark the post-rigor (\texttt{v0.1-prepaper-rigor}) and
post-airtight-matrix (\texttt{v0.2-airtight-matrix}) checkpoints.
